\documentclass[11pt]{article}
\usepackage[T1]{fontenc}
\usepackage{amsmath,amssymb,amsthm}
\usepackage[margin=1in]{geometry}
\usepackage{booktabs}
\usepackage[hidelinks,pdftitle={The Smallest Abelian Cayley Graphs with a Goemans--Linial Gap},pdfauthor={Xiaojian Zeng}]{hyperref}

\newtheorem{theorem}{Theorem}
\newtheorem{lemma}[theorem]{Lemma}
\newtheorem{proposition}[theorem]{Proposition}
\DeclareMathOperator{\Cay}{Cay}
\DeclareMathOperator{\SDP}{SDP}

\title{The Smallest Abelian Cayley Graphs with a Goemans--Linial Gap}
\author{Xiaojian Zeng\\
PBC School of Finance, Tsinghua University, 100190\\
\texttt{xiaojianzeng@foxmail.com}}
\date{}

\begin{document}
\maketitle

\begin{abstract}
For a connected simple Cayley graph of a finite Abelian group, with uniform
edge and pair demands, we determine the smallest group order at which the
Goemans--Linial relaxation is not exact.  The minimum is ten.  There are
exactly three non-exact graphs of that order, up to graph isomorphism, and
each has a single Abelian Cayley-presentation orbit.  Their exact cut and SDP
values are
\[
(4/5,3/4),\qquad (4/5,(15+\sqrt5)/22),\qquad
(5/6,(10+\sqrt5)/15).
\]
The proof combines the exponent-at-most-four theorem with a complete exact
enumeration of the residual cyclic cases.  The finite part is accompanied by
machine-checkable rational and quadratic primal--dual identities.
\end{abstract}

\section{Statement and context}

Let $\Gamma$ be a finite Abelian group and let
$S\subseteq\Gamma\setminus\{0\}$ be a set such that $S=-S$ and
$\langle S\rangle=\Gamma$.  Thus $G=\Cay(\Gamma,S)$ is connected, simple,
undirected, and $|S|$-regular.  For a nonempty proper set $A\subset\Gamma$,
put
\[
 \psi_G(A)=\frac{|\Gamma|\,|\partial A|}
 {|S|\,|A|(|\Gamma|-|A|)},\qquad
 \psi(G)=\min_A\psi_G(A).
\]
For a squared Euclidean semimetric $d$ satisfying all triangle inequalities,
write
\[
 N_G(d)=\mathbb E_{x\in\Gamma,s\in S}d(x,x+s),\qquad
 D(d)=\mathbb E_{x,y\in\Gamma}d(x,y),
\]
and define
\[
 \SDP_{\rm GL}(G)=\inf_{D(d)>0}\frac{N_G(d)}{D(d)}.
\]
Every cut metric is feasible, so $\SDP_{\rm GL}(G)\leq\psi(G)$.

Stamoulis proved exactness for groups of exponent at most four and exhibited
$\Cay(\mathbb Z_{10},\{\pm1,\pm3\})$ with SDP value $3/4$ and cut value
$4/5$ \cite[Corollary~4.7 and Section~6.4]{Stamoulis2026}.  The question left
by this ten-vertex seed is whether a smaller example exists and what the full
order-ten classification is.  Earlier work on low-degree Abelian Cayley
graphs gives a $(1+\varepsilon)$-approximation algorithm and spectral
structure rather than this exact finite classification
\cite{DorsiEtAl2025}.  The literature comparison here is deliberately
limited: no priority claim beyond these inspected sources is made.

\begin{theorem}\label{thm:main}
The least order of a finite Abelian group admitting an admissible presentation
with $\psi(G)>\SDP_{\rm GL}(G)$ is $10$.  The $16$ presentation orbits on
$\mathbb Z_{10}$ yield pairwise nonisomorphic graphs.  Exactly three are
non-exact:
\begin{center}
\begin{tabular}{c@{\qquad}c@{\qquad}c}
\toprule
$S$ & $\psi(G)$ & $\SDP_{\rm GL}(G)$\\
\midrule
$\{1,3,7,9\}$ & $4/5$ & $3/4$\\
$\{1,4,6,9\}$ & $4/5$ & $(15+\sqrt5)/22$\\
$\{1,2,5,8,9\}$ & $5/6$ & $(10+\sqrt5)/15$\\
\bottomrule
\end{tabular}
\end{center}
Each of these graph-isomorphism classes has one Cayley-presentation orbit.
\end{theorem}

\section{The cyclic linear program}

We first reduce every residual cyclic instance to a finite exact problem.

\begin{lemma}\label{lem:fourier}
For $G=\Cay(\mathbb Z_n,S)$, an optimum GL semimetric may be taken in the
form
\[
 f(t)=\sum_{k=1}^{\lfloor n/2\rfloor}b_k
 \bigl(1-\cos(2\pi kt/n)\bigr),\qquad b_k\geq0.
\]
After the normalization $D(f)=1$, one has $\sum_k b_k=1$, and feasibility is
equivalent to the finite family
\[
 f(a)+f(b)-f(a+b)\geq0\qquad(a,b\in\mathbb Z_n).
\]
Consequently the cyclic GL value is the optimum of a finite linear program.
\end{lemma}

\begin{proof}
Average a feasible metric over all simultaneous translations.  Convexity of
the squared-Euclidean cone and of the triangle inequalities preserves
feasibility, while the Cayley numerator and uniform-pair denominator are
unchanged.  A translation-invariant squared Euclidean semimetric is a
conditionally negative-type function on $\mathbb Z_n$; this is the finite
cyclic specialization of Schoenberg's negative-type correspondence
\cite{Schoenberg1938}.  Diagonalizing its
circulant matrix by the characters, and pairing conjugate characters, gives
exactly the displayed representation with nonnegative coefficients; the
converse follows by adding the corresponding squared chord metrics.  Averaging
over $t$ gives $D(f)=\sum_k b_k$.  The remaining metric axioms are precisely
the displayed triangle inequalities.  Normalization confines $b$ to a closed
subset of a simplex, so the infimum is attained.
\end{proof}

For certificates that need no negative-type information, set
$x_i=f(i)=f(n-i)$, $1\leq i\leq\lfloor n/2\rfloor$.  Let $c\cdot x=N_G(f)$,
$w\cdot x=D(f)$, and let $R_{a,b}\cdot x=f(a)+f(b)-f(a+b)$.  An identity
\begin{equation}\label{eq:rationaldual}
 c-Lw=\sum_{a,b}\lambda_{a,b}R_{a,b}+\sum_i\mu_i e_i,
 \qquad \lambda_{a,b},\mu_i\geq0,
\end{equation}
proves $N_G(f)\geq L D(f)$.  A cut attaining $L$ proves equality.  All
rational instances used below are certified by identities of
form~\eqref{eq:rationaldual}, stored as exact numerator--denominator pairs in
the accompanying verifier.

Four order-ten cases use the full Fourier program.  Put
\[
 u_k(t)=1-\cos(2\pi kt/10),\quad
 T_{a,b,k}=u_k(a)+u_k(b)-u_k(a+b),\quad
 c_k=|S|^{-1}\sum_{s\in S}u_k(s).
\]
For each row below the displayed $b^*$ is nonnegative, sums to one, satisfies
all triangle inequalities, and has objective $L$.  The exact verifier also
checks the matching identity
\begin{equation}\label{eq:quadraticdual}
 c-L\mathbf1=\sum_{a,b}\lambda_{a,b}T_{a,b}+\mu,
 \qquad \lambda_{a,b}\geq0,\quad \mu_k\geq0,
\end{equation}
in $\mathbb Q(\sqrt5)$.  Only nonzero $\lambda$ values are shown.

\begin{center}\small
\begin{tabular}{p{.23\textwidth}p{.14\textwidth}p{.50\textwidth}}
\toprule
$S$ and $L$ & nonzero $\lambda$ & $b^*$ and $\mu$\\
\midrule
$\{1,3,7,9\}$, $3/4$
& none
& $b^*=((5-\sqrt5)/10,0,(5+\sqrt5)/10,0,0)$;
  $\mu=(0,1/2,0,1/2,5/4)$\\[2pt]
$\{1,4,6,9\}$, $(15+\sqrt5)/22$
& $\lambda_{1,1}=\lambda_{4,4}=(7-\sqrt5)/88$;
  $\lambda_{1,6}=(7-\sqrt5)/44$
& $b^*=(4/11-6\sqrt5/55,1/11+3\sqrt5/11,
4/11-6\sqrt5/55,0,2/11-3\sqrt5/55)$;
$\mu=(0,0,0,(15+\sqrt5)/44,0)$\\[2pt]
$\{1,2,5,8,9\}$, $(10+\sqrt5)/15$
& $\lambda_{1,2}=(5-\sqrt5)/25$;
  $\lambda_{2,2}=(5-\sqrt5)/75$
& $b^*=(\sqrt5/3,(3-\sqrt5)/6,0,(3-\sqrt5)/6,0)$;
$\mu=(0,0,2/5,0,(8-\sqrt5)/15)$\\[2pt]
$\{1,2,3,5,7,8,9\}$, $20/21$
& $\lambda_{2,2}=2/21$
& $b^*=(1/3,0,(3-\sqrt5)/6,(1+\sqrt5)/6,0)$;
$\mu=(0,0,0,0,10/21)$\\
\bottomrule
\end{tabular}
\end{center}

\section{Exhaustion and proof of the theorem}

We isolate a small family that admits a direct proof.  Write $K_{p\times q}$
for the complete multipartite graph with $p$ equal parts of size $q$.

\begin{lemma}\label{lem:multipartite}
$\psi(K_{p\times q})=\SDP_{\rm GL}(K_{p\times q})=1$ for $p,q\geq2$.
\end{lemma}

\begin{proof}
For vectors realizing a feasible squared Euclidean semimetric, let $m_i$ and
$V_i$ be the centroid and mean squared deviation in part $i$.  The mean
squared distance within a part is $2p^{-1}\sum_iV_i$.  The mean over ordered
pairs in distinct parts equals this quantity plus
\[
 \frac1{p(p-1)}\sum_{i\ne j}\|m_i-m_j\|^2,
\]
and is therefore no smaller.  The latter mean is the graph numerator, while
the all-pairs denominator is a convex combination of the two means.  Hence
$N_G(d)\geq D(d)$.

For a cut meeting the parts in $a_i$ vertices and $a=\sum_i a_i$, direct
counting and Cauchy--Schwarz give
\[
 p|\partial A|-(p-1)a(pq-a)=p\sum_i a_i^2-a^2\geq0.
\]
Thus every cut has value at least one, and choosing one vertex from every part
attains one.
\end{proof}

We now describe the exhaustive step.  For each $5\leq n\leq10$, enumerate
every nonempty subset of $\mathbb Z_n\setminus\{0\}$, retain it exactly when
$S=-S$ and $\gcd(n,S)=1$, and quotient by multiplication by the units modulo
$n$.  The resulting orbit counts are
\[
\begin{array}{c|rrrrrr}
n&5&6&7&8&9&10\\ \hline
\text{presentation orbits}&2&5&3&8&6&16.
\end{array}
\]
For each representative, every cut/complement pair is tested exactly once by
requiring the representative side to omit $n-1$.  Boundary counts are
integers, hence this computes $\psi$ exactly.  Of the $24$ representatives
below ten, $21$ have a rational certificate
\eqref{eq:rationaldual}; the remaining three are
\[
(6,\{1,3,5\}),\quad (8,\{1,3,5,7\}),\quad
(9,\{1,2,4,5,7,8\}),
\]
namely $K_{2\times3}$, $K_{2\times4}$, and $K_{3\times3}$, covered by
Lemma~\ref{lem:multipartite}.  This style of canonical exact cut enumeration
also appears, in a different graph-reconstruction setting, in
\cite{Zeng2026Multiplicity}; here every decisive identity is recomputed in
exact arithmetic rather than accepted from numerical optimization.

Every noncyclic Abelian group of order at most nine has exponent at most four;
the same is true for cyclic groups of order at most four.  Stamoulis's
exactness theorem therefore covers them, and the preceding cyclic exhaustion
covers everything else.  Hence every admissible presentation below order ten
is exact.

At order ten, $\mathbb Z_{10}$ is the only Abelian group.  Eleven of its
$16$ presentation orbits have rational certificates
\eqref{eq:rationaldual}; the all-odd generator set is $K_{2\times5}$; and the
remaining four cases have certificates~\eqref{eq:quadraticdual}.  Exact cut
enumeration gives the values in Theorem~\ref{thm:main}, and all other cases
are exact.  The three inequalities are strict because
$5\sqrt5<13$ and $2\sqrt5<5$ (square their positive sides).

Finally, for every one of the $16$ representatives the verifier constructs
the integer adjacency matrix and computes its characteristic polynomial
exactly.  The $16$ polynomials are pairwise distinct.  Since isomorphic graphs
have the same characteristic polynomial, distinct presentation orbits cannot
collide under ordinary graph isomorphism.  This proves both the asserted
three-class list and the absence of a second inequivalent presentation for
any of those classes.  Together with exactness below ten, this completes the
proof of Theorem~\ref{thm:main}.

\section*{Computational reproducibility and assistance disclosure}

The supplementary verifier is \path{evidence/exact_cyclic_gl.py}.  It uses
integer and rational arithmetic (Python's \texttt{Fraction} class), together
with symbolic arithmetic in $\mathbb Q(\sqrt5)$ via SymPy.  It enumerates all $40$ residual cyclic
presentation orbits, every cut, every stored dual identity, all Fourier primal
slacks, and the $16$ order-ten characteristic polynomials.  It is run from the
project root with
\begin{center}
\texttt{python3}\ \path{evidence/exact_cyclic_gl.py}.
\end{center}
No floating-point SDP value is used as proof.

An OpenAI Codex model assisted with organizing the accepted proof dossier,
checking bibliographic records, and preparing the LaTeX exposition.  The
finite certificates were independently recomputed by the exact verifier, and
the author takes responsibility for all statements and calculations.

\bibliographystyle{plain}
\bibliography{references}

\end{document}
